\documentclass[11pt,a4paper]{article}
\usepackage[margin=0.85in,top=1in,bottom=1in]{geometry}
\usepackage{amsmath,amssymb,amsfonts}
\usepackage{graphicx}
\usepackage{float}
\usepackage{enumitem}
\usepackage{microtype}
\usepackage{caption}
\usepackage[hidelinks]{hyperref}
\usepackage[most,breakable]{tcolorbox}
\tcbuselibrary{breakable,skins}
\usepackage{fontspec}
\setmainfont{DejaVu Serif}
\setsansfont{DejaVu Sans}
\setmonofont{DejaVu Sans Mono}
\captionsetup{font=small,labelfont=bf,skip=4pt}
\setlist{leftmargin=1.5em,topsep=2pt}

%% Breakable problem card — prevents content cutoff across pages
\newtcolorbox{solcard}[3]{
  breakable, enhanced,
  colback=white, colframe=gray!50, arc=2pt,
  title={\textbf{#1}\hfill\textsf{\small #3\,pts}},
  fonttitle=\normalsize\sffamily\bfseries,
  before upper={\small\textit{#2}\par\smallskip\hrule height 0.4pt\smallskip},
  top=4pt, bottom=6pt, left=7pt, right=7pt,
  parbox=false,
}

\title{\textbf{Metronome Sync --- Solution}}
\author{\normalsize Y22 (2022) — English translation}
\date{}
\begin{document}\maketitle\vspace{-0.5em}
\begin{solcard}{A1}{Using the metronome oscillation graph, find with an accuracy of 10\%: the angular frequency $\omega$, the steady-state amplitude $A$, the damping coefficient $\gamma$ and the angular velocity $\omega_D$ reported by the trigger mechanism upon impact. It is known that $\frac{\omega_D}\omega,\ \gamma\ll1$.}{1.50}

It is clear from the graph that by $12$ the change in amplitude stops. The steady-state amplitude is $A=0.3$ rad. Let us immediately note that at such angles the approximation error $\sin x \approx x$ does not exceed $1.5\%$. In particular, for the purposes of this paragraph, frequency can be assumed to be independent of amplitude. The frequency can be found using the series method, for example $$ \omega = 2\pi \frac{16}{14.5c - 9.5c} = 20\; \text{s}^{-1}. $$Колебания в остутствии подкачки энергии описываются equationм $$ ml^2\ddot{\theta} = -2ml^2\gamma\dot{\theta} - mgl \sin \theta $$    (angular momentum theorem for pendulumа). Заменяя в силу сказаного выше $\sin \theta$ на $\theta$, $g$ на $l\omega^2$ and сокращая на $ml^2$ we get equation гармонических колебаний: $$ \ddot{\theta} + 2\gamma \dot{\theta} + \omega^2 \theta = 0. $$    Его solution с начальными условиями $\theta(0) = 0$, $\dot{\theta}(0) = A\omega$ есть $$ \theta(t) = A e^{-\gamma t} \sin \omega t. $$    Denote $A_n$ value $A$ на $n$-том полуperiodе. Каждый полуperiod $\dot {\theta}$ увеличивается на $\omega_D$, therefore $$ A_{n+1} = A_n e^{-\frac{\gamma \pi}{\omega}} + \frac{\omega_D}{\omega}. $$    Limitьная amplitude находится from the condition $$ A = A e^{-\frac{\gamma \pi}{\omega}} + \frac{\omega_D}{\omega} $$    hence $A = \frac{\omega_D}{\omega} \left(1 - e^{-\frac{\gamma \pi}{\omega}}\right)^{-1}$. Рекуррентное соratio на $A_n$ можно переписать так: $$ A - A_{n+1} = (A - A_n) e^{-\frac{\gamma \pi}{\omega}}. $$ Это позволяет найти $\gamma$ по нескольким пикам (в предположении что $\gamma$ мало and максимальное value $\theta$ equals $A$): $$ \gamma = \frac{\omega}{\pi} \frac{1}{m - n} \ln \frac{A-A_n}{A-A_m}. $$    For example, $A_n = 0.14$, $A_m = 0.19$, $m - n = 6$ (число полуperiodов вдвое больше числа periodов). Поrayается $\gamma = 0.40$. $\gamma/\omega = 0.02$ and затухание действительно мало. Остается найти $\omega_D$. $$ \omega_D = \omega A \left(1 - e^{-\frac{\gamma \pi}{\omega}}\right) = 0.37 \; \frac{\mathrm{rad}}с. $$

\par\smallskip\noindent\textbf{Answer:} \textbackslash{}textbf{Answer:} \[\omega = 20~\text{s}^{-1},\\A=0.3~\mathrm{rad},\\\gamma = 0.40~\mathrm{s}^{-1},\\\omega_D = 0.37~\frac{\mathrm{rad}}{\mathrm{s}}\]
\end{solcard}

\begin{solcard}{B1}{Express $\dot{x}(t)$ in terms of $A(t)$ and $A^*(t)$.}{0.20}

The answer is obtained by direct differentiation.

\par\smallskip\noindent\textbf{Answer:} \textbackslash{}textbf{Answer:} \[\dot{x}(t) = \frac{1}{2}(\dot{A}e^{i\omega t}+i\omega A e^{i\omega t} + \dot{A^*}e^{-i\omega t}-i\omega A^* e^{-i\omega t})\]
\end{solcard}

\begin{solcard}{B2}{Express $\ddot{x}(t)$ in terms of $A(t)$, $A^*(t)$, $\dot{A}(t)$.}{0.30}

$$ \ddot{x}(t) = -\frac{\omega^2}{2}\left(A e^{i\omega t} + A^*e^{-i\omega t}\right) + \frac{1}{2}\left( i \omega \dot{A} e^{i\omega t} - i \omega A^* e^{-i\omega t} \right) $$Taking into account that мы наложили condition $\dot{A} e^{i \omega t} + \dot{A^*} e^{-i \omega t} = 0$, we simplify this expression and get the answer.

\par\smallskip\noindent\textbf{Answer:} \textbackslash{}textbf{Answer:} \[\ddot{x}=-\frac{1}{2}\omega^2(A e^{i \omega t} + A^*e^{-i\omega t}) + i \omega \dot{A} e^{i \omega t}\]
\end{solcard}

\begin{solcard}{B3}{Using the vibration equation, find an expression for $\dot{A}$ in terms of $A$, $\gamma$, $\omega$, $\varepsilon$.}{0.50}

We substitute the result into the oscillation equation: $$ -\frac{\omega^2}{2}\left(A e^{i\omega t} + A^*e^{-i\omega t}\right) + i \omega \dot{A} e^{i\omega t} + \gamma \left(i \omega A e^{i \omega t} -i\omega A^*e^{-i\omega t}\right) + \frac{\omega^2}{2}\left(A e^{i\omega t} + A^*e^{-i\omega t}\right) =\\= \varepsilon \frac{1}{8} \left(A e^{i \omega t} + A^*e^{-i\omega t}\right)^3 . $$    Hence we express $\dot{A}$.

\par\smallskip\noindent\textbf{Answer:} \textbackslash{}textbf{Answer:} \[\dot{A} = \frac{1}{2i\omega}e^{-i \omega t}\left(-2\gamma(i \omega A e^{i\omega t} - i \omega A^* e^{-i\omega t}) +\\+ \frac{\varepsilon}{4}(A^3 e^{3i\omega t} + 3A |A|^2 e^{i\omega t} + A^{*3} e^{-3i\omega t} + 3A^* |A|^2 e^{-i\omega t})\right)\]
\end{solcard}

\begin{solcard}{B4}{Given this, write a simplified equation for $\dot{A}(t)$.}{0.20}

In the equation obtained in $\bf B3$, you simply need to discard all terms containing $e^{-in\omega t}$, $n\neq0$.

\par\smallskip\noindent\textbf{Answer:} \textbackslash{}textbf{Answer:} \[\dot{A}=-A\left(\gamma +\frac{3\varepsilon}{8\omega}i|A|^2\right)\]
\end{solcard}

\begin{solcard}{B5}{Let us assume that there is no friction, $\gamma = 0$. Let the initial value of $A$ be $A_0$. Solve the equation from $\bf B4$ and get $A(t)$ and $x(t)$.}{0.60}

Putting $\gamma=0$ in the equation obtained in $\bf B4$, we obtain \[\dot A=-i\frac{3\varepsilon|A|^2}{8\omega}A.\]From geometric considerations it is clear that this equation describes the rotation of the vector $A$ in the complex plane without changing its modulus. Therefore, we look for a solution to the equation in the form \[A(t)=A_0e^{ist}.\]We obtain \[s=-\frac{3\varepsilon|A|^2}{8\omega},\]from which we immediately obtain the answers for $A(t)$ and $x(t)$.

\par\smallskip\noindent\textbf{Answer:} \textbackslash{}textbf{Answer:} \[A(t)=A_0\exp\left[-i\frac{3\varepsilon|A_0|^2}{8\omega}t\right]\\x(t)=\operatorname{Re}\left\{A_0\exp\left[i\left(\omega-\frac{3\varepsilon|A_0|^2}{8\omega}\right)t\right]+A^*_0\exp\left[-i\left(\omega-\frac{3\varepsilon|A_0|^2}{8\omega}\right)t\right]\right\}\]
\end{solcard}

\begin{solcard}{B6}{Find the first correction $\Delta \omega$ to the frequency of oscillations described by the equation \[ \ddot{x} + \omega^2 \sin x = 0 \] with a small amplitude $A \ll 1$. What is $\frac{\Delta \omega}{\omega}$ equal to with the amplitude found in $\bf A1$?}{0.40}

Let us expand the sine into a Taylor series:\[\sin x\approx x-\frac16x^3.\]Thus, to a first approximation, the equation of oscillation of a pendulum is described by the equation from $\bf B4$ with\[\gamma=0,\quad\varepsilon=\frac16\omega^2.\]From $\bf B5$ we know that its solution can be represented in the form\[A(t) = A_0 \exp \left(-i \frac{|A_0|^2 \omega t}{16}\right)\implies\\\implies x(t)=\operatorname{Re}\left\{A_0\exp\left[i\left(1-\frac{1}{16}A^2\right)\omega t\right]\right\}.\]From here we immediately get the answer.

\par\smallskip\noindent\textbf{Answer:} \textbackslash{}textbf{Answer:} \[\Delta\omega=-\frac{1}{16}A^2\omega,\\\frac{\Delta\omega}\omega\approx -6\cdot 10^{-3}\]
\end{solcard}

\begin{solcard}{C1}{Let's start with an auxiliary task. Let the metronomes stand on a stationary board and oscillate with known dependencies $\theta_1(t)$ and $\theta_2(t)$. In first order in terms of vibration amplitudes, find the force acting on the board. Express your answer in terms of $\theta_1$, $\theta_2$, $D_1$, $D_2$ and the constants that characterize metronomes.}{0.70}

The board is acted upon by the tension forces of the pendulum threads, equal to a first approximation $mg\sin\theta_1$ and $mg\sin\theta_2$. According to Newton's third law, during a push, a force $-ml(D_1+D_2)$ also acts on the board. As a result, we get the answer.

\par\smallskip\noindent\textbf{Answer:} \textbackslash{}textbf{Answer:} \[F=mg(\sin\theta_1+\sin\theta_2)-ml(D_1+D_2)\]
\end{solcard}

\begin{solcard}{С2}{Write the equations of motion of the loads in the reference frame of the board, that is, express $\ddot{\theta}_1$, $\ddot{\theta}_2$ in terms of $\theta_1$, $\theta_2$, $\dot{\theta}_1$, $\dot{\theta}_2$, $\omega$, $\gamma$, $\alpha$, $D_1$, $D_2$.}{1.00}

The reference system of the board is a non-inertial reference system moving with acceleration $a=\frac FM$, then in Newton’s law for loads, as a first approximation, an additional term $-ma$ should appear. Rewriting it through the angles $\theta_1$ and $\theta_2$, we get the answer.

\par\smallskip\noindent\textbf{Answer:} \textbackslash{}textbf{Answer:} \[\begin{cases}\ddot\theta_1+2\gamma\dot\theta_1+\omega^2\theta_1=\alpha\left[-(\theta_1+\theta_2)\omega^2+(D_1+D_2)\right]\\\ddot\theta_2+2\gamma\dot\theta_2+\omega^2\theta_2=\alpha\left[-(\theta_1+\theta_2)\omega^2+(D_1+D_2)\right]\end{cases}\]
\end{solcard}

\begin{solcard}{C3}{The phase diagram of oscillation is an image of the movement of a pendulum on a plane with coordinates $\dot{\theta}$, $\omega \theta$. Let us assume that an oscillation mode has been established in which their amplitudes are constant, but at the same time there is a certain phase difference $\psi$. Construct the phase diagram of the first oscillator for one period. Indicate on it the points at which the metronome triggers are triggered. Find the change in the phase difference between the metronomes over the period $\Delta \psi$.}{2.00}

As found in $\bf A1$, the moment the axis $\dot{\theta}$ is crossed, the value of $\dot{\theta}$ changes to $\gamma\pi A$. When after some time this axis is crossed by the first pendulum, the second one also experiences a disturbance and its $\dot{\theta}$ changes to $\alpha \gamma \pi A$. The phase of oscillations, up to a constant equal to the angle through which the point on the phase portrait has rotated, decreases by \[\frac{\alpha \gamma \pi A \sin \psi}{\omega A} = \frac{\alpha \gamma\pi}{\omega}\sin\psi.\]When in the same way the click of the second metronome affects the phase of the first, the phase difference between them decreases again. Thus, over a half-period \[\Delta \psi = - 2 \frac{\alpha \gamma \pi}{\omega}\sin \psi.\]Over a period, this value will be twice as large.

\begin{figure}[H]
  \centering
  \includegraphics[width=0.72\textwidth]{T4_S_EN_assets/fig1.jpg}
\end{figure}
\end{solcard}

\begin{solcard}{С4}{Show that the change in phase difference between the oscillators satisfies the Kuramoto equation$$ \dot{\psi} = - B \sin \psi. $$Find expression for постоянной $B$ and ее численное value, if $\alpha=10^{-2}$.}{0.70}

Since we consider the attenuation to be small, the change in phases over a period turns out to be small. Therefore, we can go to the continuous limit and get the answer.

\par\smallskip\noindent\textbf{Answer:} \textbackslash{}textbf{Answer:} \[B=2\gamma\alpha= 0.8 \cdot 10^{-3}~\mathrm{s}^{-1}\]
\end{solcard}

\begin{solcard}{С5}{Let us assume that the natural frequencies of the oscillators are slightly different and are equal to $\omega_1$ and $\omega_2$, where $\Delta \omega = \omega_1 - \omega_2 \ll \omega_1$. Find an expression for $\dot{\psi}$ in this case. At what maximum frequency difference $\Delta \omega_{max}$ is synchronization possible?}{0.80}

Let us rewrite the equations of motion of pendulums from $\bf C2$, taking into account the different frequencies of the pendulums. We get $$ \begin{cases}\ddot{\theta}_1+2\gamma\dot{\theta}_1+\omega_1^2 \theta_1 = \alpha\left[-(\theta_1+\theta_2)\omega^2+(D_1+D_2)\right]\\\ddot{\theta}_2+2\gamma\dot{\theta}_2+\omega_2^2 \theta_2 = \alpha\left[-(\theta_1+\theta_2)\omega^2+(D_1+D_2)\right]\end{cases} $$Then similarly⟦0⟧Синхронизация возможна, когда это equation имеет корни, hence ссу же поrayается answer for $\Delta\omega_\max$.

\par\smallskip\noindent\textbf{Answer:} \textbackslash{}textbf{Answer:} \[\dot\psi=\omega_1-\omega_2-\frac{3\omega^2A^3\alpha}{16\gamma}\sin\psi,\\\Delta\omega_\max=\frac{3\omega^3A^3\alpha}{16\gamma}\]
\end{solcard}

\begin{solcard}{С6}{Let the frequency difference $\Delta \omega$ be less than the maximum. Find the steady-state phase difference $\psi$ of the oscillators.}{0.20}

With a steady phase difference between the metronomes $\dot\psi=0$. The answer is obtained by substituting this into the expression obtained in $\bf C5$.

\par\smallskip\noindent\textbf{Answer:} \textbackslash{}textbf{Answer:} \[\psi=\arcsin\frac{16\gamma\Delta\omega}{3\omega^3A^3\alpha}\]
\end{solcard}

\begin{solcard}{C7}{Let the frequency difference in a system of two metronomes be $\Delta \omega =0.1 c^{-1}$. In this case, the phase difference is $\psi = 0.5 ~\text{rad}$.  At some point in time, the weight of one of the metronomes shifts slightly and the frequency difference becomes $0$. Find the time $t$ during which the metronomes are completely synchronized, that is, their phase difference becomes less than $0.01~\text{rad}$.}{0.90}

It is necessary to solve the Kuramoto equation with the initial condition\[\psi(0)=0.5.\]We have\[\dot\psi=-B\sin\psi\implies\int_{0.5}^{\psi(t)}\frac{\mathrm d\psi}{\sin\psi}=\ln\bigg|\operatorname{tg}\frac\psi2\bigg|\Bigg|_{0.5}^{\psi(t)}=-Bt\implies\\\implies t(\psi)=\frac{\sin\psi_0}{\Delta\omega}\ln\left|\frac{\operatorname{tg}\frac\psi2}{\operatorname{tg}0.05}\right|.\]

\par\smallskip\noindent\textbf{Answer:} \textbackslash{}textbf{Answer:} \[t=\frac{\sin\psi}{\Delta\omega}\ln\left|\frac{\operatorname{tg}\frac\psi2}{\operatorname{tg}0.05}\right|\approx18.9~\mathrm{s}\]
\end{solcard}

\end{document}